\subsection{All $A_5$ chief extensions over finite odd-order groups}

\begin{statement}{D.29}{Proved}{Exact peripheral count invariance for every $A_5$ chief step over an odd-order group}
Let
\[
 1\longrightarrow A\longrightarrow E\xrightarrow{\pi}D
 \longrightarrow1,
 \qquad A\simeq A_5^{\Omega},
 \tag{D.29.1}
\]
be a finite extension, with $D$ of odd order and $A$ a
nonabelian chief factor of $E$.  Let
$q_E:\widehat F\twoheadrightarrow E$ be an epimorphism,
$X_i=q_E(x_i)$, $q=\pi q_E$, and $d_i=q(x_i)$.
Given $\alpha\in\operatorname{Aut}(D)$,
$\lambda\in\widehat\Z^\times$, and specified elements
$h_i\in D$ such that
\[
           \alpha(d_i)=h_i d_i^\lambda h_i^{-1},
 \tag{D.29.2}
\]
choose any lifts $H_i\in E$ of $h_i$.
There is an epimorphism $f:\widehat F\twoheadrightarrow E$
such that
\[
 \pi f=\alpha q,
 \qquad
 f(x_i)=a_iH_iX_i^\lambda H_i^{-1}a_i^{-1}
              \quad\text{for some }a_i\in A.
 \tag{D.29.3}
\]
The number of these literal epimorphisms is exactly the
number for $(\alpha,\lambda,h_1,h_2,h_3)=(1,1,1,1,1)$.
In particular it is positive.  There is no degree or
fixed-point restriction on the transitive factor action
on $\Omega$.

The same assertion holds in every actual fibre product
$E\times_D L$, for a finite odd-order group $L$ and a prescribed
peripheral automorphism of $L$ inducing $\alpha$ on $D$:
there is a generating lift over that automorphism, with
its specified peripheral conjugators retained in the
$L$-coordinate.  Thus every $A_5$ chief extension over a
finite odd-order group, including its coupled pullbacks, admits
all the peripheral shadows described in
\textup{(D.29.2)--(D.29.3)}.
In particular this applies over every finite $\ell$-group
for every odd prime $\ell$.
\end{statement}

\begin{proof}
Put $S=A_5$.  We first give the exact map from
\textup{(D.29.1)} to a permutation extension, including
the possible outer twisting.  Every automorphism of
$S^{\Omega}$ permutes its simple direct factors and
acts on each factor by an automorphism of $S$.
The computation in D.28 gives
\[
 Z(S)=1,\qquad
 \operatorname{Out}(A)=\mathbf F_2^{\Omega}
                         \rtimes\operatorname{Sym}(\Omega).
 \tag{D.29.4}
\]
The factor action $\rho:D\to\operatorname{Sym}(\Omega)$
is transitive: a union of factor orbits would otherwise
give a proper nonzero normal subgroup inside the chief
kernel.  The outer action of $D$ is of the form
$(b_d,\rho(d))$, with
$b_{de}=b_d+\rho(d)b_e$.
Since $|D|$ is odd, the vector
\[
 B=|D|^{-1}\sum_{e\in D}b_e\in\mathbf F_2^{\Omega}
 \quad\text{satisfies}\quad b_d=B-\rho(d)B.
 \tag{D.29.5}
\]
This follows by summing the cocycle identity over $e$.
Choose a componentwise automorphism $\tau_B$ of $A$
with outer class $B$.  Such a representative is explicit:
on each component with value one use conjugation by the
odd permutation $(01)\in S_5$; on the other components
use the identity.  Conjugating the outer action by $B$
therefore makes it exactly $\rho$.

Here is the corresponding extension isomorphism, rather
than only an identification of outer actions.  Since
$Z(A)=1$, the map
\[
 E\longrightarrow\operatorname{Aut}(A)
                       \times_{\operatorname{Out}(A)}D,
       \quad e\longmapsto(\operatorname{Ad}_e,\pi(e))
 \tag{D.29.6}
\]
is an isomorphism.  It is injective because its kernel
lies in $A\cap C_E(A)=Z(A)$.  For surjectivity, an
automorphism in the required outer class differs from
the conjugation of any chosen lift of $d$ by a unique
inner automorphism of $A$.
After applying $\tau_B^{-1}$ to this kernel identification,
there is for every $e\in E$ a unique $a_e\in A$ with
\[
 \tau_B^{-1}\operatorname{Ad}_e\tau_B
              =\operatorname{Inn}(a_e)\rho(\pi(e)).
 \tag{D.29.7}
\]
Uniqueness uses centerlessness.  Multiplication gives
$a_{ef}=a_e\rho(\pi(e))(a_f)$, so
$e\mapsto(a_e,\pi(e))$ is the actual isomorphism
$E\simeq A\rtimes_\rho D$.  For example, after choosing
a lift $e_d$ of $d$, the inverse image of $(a,d)$ is
$\tau_B(a a_{e_d}^{-1})e_d$.
This map transports the specified source elements and
all conjugators in \textup{(D.29.3)}.  We may perform
the rest of the calculation in this displayed
semidirect product.

Write the factor action on the right as in D.28,
choose $\omega_0\in\Omega$, and let $D_*$ be its
stabilizer.  Thus $\Omega\simeq D_*\backslash D$,
and put $n=|\Omega|$.  For every subgroup
$D_*\leq C\leq D$ use the two actual covers
\[
 H_C=q^{-1}(C),\qquad
 H'_C=(\alpha q)^{-1}(C),\qquad
 K=\ker q=\ker(\alpha q).
 \tag{D.29.8}
\]
The action on $\Omega$ is allowed to be nonfaithful;
the quotient $H_C/K=C$ still has odd order.

The original peripheral elements $X_i$ determine
conjugacy classes of $S$ on every induced boundary
component of $H_{D_*}$.  These are their cycle-product
classes under the exact map \textup{(D.28.15)}.
The prescribed elements
$H_iX_i^\lambda H_i^{-1}$ determine the new boundary
classes for $H'_{D_*}$.
For each $i$ and $C$, the boundary comparison is the
explicit bijection
\[
 \begin{split}
 \sigma_{i,C}:
 C\backslash D/\langle\alpha(d_i)\rangle
       &\longrightarrow C\backslash D/\langle d_i\rangle,\\
       C d\langle\alpha(d_i)\rangle
       &\longmapsto C d h_i\langle d_i\rangle.
 \end{split}
 \tag{D.29.9}
\]
The identity $\alpha(d_i)h_i=h_i d_i^\lambda$
proves that this is well defined, and that it preserves
cycle lengths.  The common exponent is a unit, so
$d_i$ and $d_i^\lambda$ have the same cycles.
Right multiplication by $h_i$ also commutes with
every projection $D_*\backslash D\to C\backslash D$.
Consequently \textup{(D.29.9)} respects the partition
of boundary components under each covering
$H_{D_*}\leq H_C$ and preserves every relative
boundary degree.

The new class at a boundary sent to an old class
$C_\gamma$ by \textup{(D.29.9)} is exactly
$C_\gamma^\lambda$.  Indeed a cycle of length $e$
in $X_i^\lambda$ has product equal to the corresponding
coordinate of $X_i^{\lambda e}$, which is the
$\lambda$th power of the cycle product of $X_i$.
This remains true when $d_i^e\ne1$ in $D$: it fixes
that coordinate, which is all that is needed for the
power identity.  Conjugation by the kernel component
of $H_i$ changes only the conjugating element in $S$.
This proves the class comparison with the specified
base conjugators, including fixed peripheral points.

For each boundary component $\delta$ of $H_C$, let
the components $\gamma$ above it in $H_{D_*}$ have
relative degrees $r_\gamma$.  Define the actual
conjugacy-invariant subset of $S$
\[
 T_{C,\delta}=\bigcap_{\gamma\mapsto\delta}
          \{s\in S:s^{r_\gamma}\in C_\gamma\}.
 \tag{D.29.10}
\]
A map $\psi:H_C\to S$ restricts to one with the
specified classes on $H_{D_*}$ exactly when its
$\delta$-boundary value belongs to $T_{C,\delta}$.
This uses the actual covering identity that a
boundary generator upstairs is conjugate, inside
$H_C$, to the $r_\gamma$th power of the generator
downstairs.  No condition on an unspecified path
value remains because $C_\gamma$ is a conjugacy class.
The transported subset for $H'_C$ is
\[
       T'_{C,\delta'}=
       T_{C,\sigma_{i,C}(\delta')}^\lambda.
 \tag{D.29.11}
\]
The unit-power map on a finite group is a bijection
and commutes with every power $r_\gamma$, which
proves this equality even when that latter power
map is not injective.

For $A_5$ all the powers in \textup{(D.29.11)} are
realized by one actual target automorphism.  Namely
let $\tau_\lambda$ be the identity for
$\lambda\equiv1,4\pmod5$ and conjugation by $(01)$
in $S_5$ for $\lambda\equiv2,3\pmod5$.
Then
\[
       \tau_\lambda(T)=T^\lambda
       \quad\text{for every conjugacy-invariant }T\subseteq A_5.
 \tag{D.29.12}
\]
The identity, double-transposition and three-cycle
classes are each preserved by every unit power.
The two five-cycle classes are exchanged precisely
by the nonsquare residues modulo five, as they are
by the indicated odd conjugation.  All eight units
modulo $30=\exp(A_5)$ are checked in the finite
certificate below.

The two covers in \textup{(D.29.8)} have the same
degree and the same boundary cycle lengths under
\textup{(D.29.9)}, hence the same genus and number
of boundary components.  Their boundary permutation
is realized by an explicit surface-group isomorphism.
In the presentation \textup{(D.28.7)}, the adjacent
boundary transformation
\[
 (c_j,c_{j+1})\longmapsto
       (c_{j+1},c_{j+1}^{-1}c_jc_{j+1})
 \tag{D.29.13}
\]
fixes all handles and other boundaries and preserves
the boundary product literally.  Its inverse sends
$(c_j,c_{j+1})$ to
$(c_jc_{j+1}c_j^{-1},c_j)$.
These actual group automorphisms realize every
permutation of the oriented boundary conjugacy
classes; the same formulas give the inverse maps
on the profinite completions.
Choosing matching surface presentations and composing
these maps yields an isomorphism
$u_C:H'_C\to H_C$ with the boundary correspondence
\textup{(D.29.9)}.
Thus
\[
                   \psi\longmapsto
                \tau_\lambda\circ\psi\circ u_C
 \tag{D.29.14}
\]
is an actual bijection between the epimorphisms
$H_C\twoheadrightarrow S$ with the conditions
\textup{(D.29.10)} and those on $H'_C$ with the
conditions \textup{(D.29.11)}.
Let their common finite cardinality be $N_C$.
These isomorphisms are chosen separately for each
$C$; compatibility among them is not assumed or
needed for the count below.

We now impose full generation of the simple-factor
kernel.  For any epimorphism
$\varphi:H_{D_*}\twoheadrightarrow S$, its
restriction to $K$ is onto.  Indeed $\varphi(K)$ is
normal in the simple group $S$, so it is either $1$
or $S$.  The former alternative would make the
even-order group $S$ a quotient of the odd-order
group $H_{D_*}/K$, which is impossible.
Put $N=\ker(\varphi|_K)$.
For $D_*\leq C\leq D$, the condition
$H_C\leq N_{\widehat F}(N)$ is equivalent to
extension of $\varphi$ to $H_C$, and the extension
is unique.  Indeed conjugation on $K/N\simeq S$
gives $H_C\to\operatorname{Aut}(S)$; its outer
part factors through the odd-order group $C$ and is
trivial because $\operatorname{Out}(S)=C_2$.
The inverse of the injective map
$S\xrightarrow{\operatorname{Inn}}\operatorname{Inn}(S)$
is the required homomorphism to $S$.
It agrees with $\varphi$ on $H_{D_*}$ by
centerlessness, and is uniquely determined by
conjugation on the surjective image of $K$.

Conversely an epimorphism $H_C\to S$ restricts
surjectively to $H_{D_*}$, because its restriction
to $K$ is already onto by the same odd-order
quotient argument.  Its boundary conditions on
that restriction are exactly
\textup{(D.29.10)}.  Hence $N_C$ is also the
number of the prescribed epimorphisms from
$H_{D_*}$ whose coordinate-kernel stabilizer
contains $H_C$.
Let $e_B$ count those with stabilizer exactly
$H_B$, for $D_*\leq B\leq D$.
On this finite subgroup poset we have
\[
 N_C=\sum_{B\geq C}e_B,\qquad
 e_{D_*}=\sum_{C\geq D_*}\mu(D_*,C)N_C,
 \tag{D.29.15}
\]
where $\mu$ is the inverse of the incidence matrix
of inclusion.  Formula \textup{(D.29.14)} proves
equality of every $N_C$ for the original and the
new problem, so \textup{(D.29.15)} proves equality
of their $e_{D_*}$ as well.

The original source epimorphism $q_E$ gives
$e_{D_*}>0$.  In fact the explicit nonabelian
Shapiro map of \textup{(D.28.13)} identifies
the coordinate maps on $K$ with
$k\mapsto\varphi(t_\omega k t_\omega^{-1})$.
Their kernels are pairwise distinct precisely
when the stabilizer of $N$ is $H_{D_*}$.
For the original epimorphism their joint image
is all of $S^\Omega$, so these kernels are
distinct.  Conversely, for the new problem
distinct kernels and surjectivity of the
coordinate maps force the full product by
the subdirect argument of D.25.
The Shapiro formula therefore turns each
new map counted by $e_{D_*}$ into an
epimorphism in \textup{(D.29.3)}, including
the prescribed $h_i$.

The literal count can also be retained exactly.
For fixed representatives $t_\omega$, every
homomorphism to $S^\Omega\rtimes D$ over the
base map has a unique kernel conjugation by
$v\in S^\Omega$ with $v(\omega_0)=1$
that makes its cocycle satisfy
$c(t_\omega)(\omega_0)=1$.
Explicitly take
$v(\omega)=c(t_\omega)(\omega_0)$ and use
\[
 c^v(g)(\omega)=v(\omega)c(g)(\omega)
                         v(\omega q(g))^{-1}.
 \tag{D.29.16}
\]
Here $q$ denotes the base map in the case being
counted; in the new case it is replaced by $\alpha q$.
The resulting cocycle is exactly the Shapiro
cocycle of $\varphi=c|_{H_{D_*}}(\omega_0)$.
Thus each surface map has exactly $60^{n-1}$
literal cocycles.  Both generation and the
specified kernel-conjugacy conditions are
preserved by these conjugations, so the
literal epimorphism count is
\[
                 60^{n-1}
                 \sum_{C\geq D_*}\mu(D_*,C)N_C.
 \tag{D.29.17}
\]
Its proved invariance gives the asserted exact
count, as well as existence.

For the coupled assertion, choose the specified
$L$-conjugators of its peripheral automorphism
and use their images in $D$ as the $h_i$ in
\textup{(D.29.2)}.  The lift just constructed
pairs with that automorphism of $L$, and
the peripheral conjugators are the literal
pairs of the chosen lifted conjugators.
For surjectivity, let $K_L$ be the kernel of the
map from $\widehat F$ to $L$.  Its image in $A$
is normal in $A$, and the quotient of $A$ by this
image is a quotient of the odd-order group
$\ker(L\to D)$.  Every normal subgroup of a direct
power of the nonabelian simple group $S$ is a
product of its simple factors: commuting such a
subgroup with each factor first includes every
factor on which its projection is nontrivial.
Thus every nontrivial quotient of $A$ has a
factor of order $60$ and has even order.
The quotient in question is therefore trivial,
which proves surjectivity of the paired map.
When $L$ is an actual $\ell$-group, its coordinate
projection is the maximal pro-$\ell$ quotient:
the perfect kernel $A$ has no nontrivial finite
$\ell$-group quotient.  These arguments use
neither a solvability theorem for odd-order
groups nor an inference from perfectness to
absence of all odd-order quotients.

The program \texttt{a\_peripheral\_a5\_invariance.py}
checks the actual outer involution of $A_5$,
all eight unit powers modulo thirty, all
thirty-two conjugacy-invariant subsets, and
their inverse images under every power modulo thirty.
It verifies \textup{(D.29.13)} and its inverse
as literal free-word identities, and checks
the averaging identity on a basis of the
outer permutation module.
It also treats the degree-five regular cover
with $x\mapsto t,y\mapsto1$, which has seven
boundaries and genus zero.  For a generating
triple of class orders $(3,5,5)$ it checks
all eight multipliers and all five possible
$y$-conjugator cosets.  In every one of the
forty cases there are sixty proper surface
maps; the corresponding literal count is
$60^5$.  The script evaluates their wreath
cycle products and the coordinate Schreier
maps, verifying surjectivity and distinct
kernels.  The code, executed output and
JSON certificate end with
\texttt{PASS\_ALL\_DEGREE\_A5\_COUNT\_COMPARISONS}.
\end{proof}

D.29 proves the complete $A_5$-chief comparison
over finite odd-order bases.  It does not
iterate automatically over an arbitrary
mixed finite base: after adjoining the
first $A_5$ factor, the base need no longer
have odd order.  The specific uses of
that hypothesis are the outer-cocycle
averaging in \textup{(D.29.5)} and the
unique inner extension in
\textup{(D.29.15)}.  Those exact maps
identify the remaining step in extending
this family toward the full cofinal
problem in D.1 and D.2.
